\chapter{Kesimpulan dan Saran}\label{cha:conclusion}

\section{Kesimpulan}
Berdasarkan hasil penelitian dan pembahasan mengenai evaluasi pengalaman pengguna pada platform \gls{digilib} IIB Darmajaya menggunakan metode \gls{ueq}, dapat ditarik beberapa kesimpulan sebagai berikut:
\begin{enumerate}
    \item Kualitas pengalaman pengguna (\gls{ux}) pada \gls{digilib} IIB Darmajaya secara keseluruhan masih berada di bawah rata-rata (\textit{below average}) untuk aspek \textit{Dependability} (0,89), \textit{Attractiveness} (1,01), dan \textit{Perspicuity} (1,02). Hal ini menunjukkan bahwa antarmuka sistem masih dirasa membingungkan, kurang menarik secara visual, alur pencarian berbelit-belit, dan filter pencarian kurang presisi.
    \item Beberapa aspek fungsionalitas dan motivasi telah dinilai di atas rata-rata (\textit{above average}), yaitu aspek \textit{Efficiency} (1,13), \textit{Stimulation} (1,04), dan \textit{Novelty} (1,00). Temuan ini menunjukkan bahwa sistem memiliki kecepatan respons yang baik dan konsep fungsional yang cukup inovatif serta relevan dengan kebutuhan mahasiswa.
    \item Hasil dari kuesioner kuantitatif didukung oleh temuan kualitatif melalui wawancara, di mana pengguna mengeluhkan antarmuka yang monoton, tata letak yang terlalu padat, serta filter pencarian yang sering kali menampilkan dokumen tidak relevan.
\end{enumerate}

\section{Saran}
Berdasarkan kesimpulan di atas, beberapa saran konstruktif yang diusulkan untuk pengembangan platform \gls{digilib} IIB Darmajaya di masa depan adalah:
\begin{enumerate}
    \item Menerapkan pendekatan \textit{User-Centered Design} (UCD) pada pengembangan bagian depan (\textit{front-end}) aplikasi untuk menyederhanakan alur navigasi (\textit{user flow}) dan mengurangi beban kognitif pengguna.
    \item Melakukan desain ulang tata letak antarmuka agar lebih modern, tidak terlalu padat, dan mengelompokkan kategori referensi secara lebih jelas dan terstruktur.
    \item Memperbaiki algoritma mesin pencari dan sistem penyaringan (filter) dokumen agar hasil pencarian literatur menjadi lebih akurat dan relevan dengan kata kunci yang dimasukkan pengguna.
\end{enumerate}